%======================================================================
\section{The Translation}\label{sec:translation}

\AP Fix a formula $\alpha$. We assign a \intro{polarity}, \intro{positive}
or \intro{negative}, to each occurrence in $\alpha$ as follows. The root
$\rt$ is \kl{positive}. If
$\subf{\alpha}{u}=\subf{\alpha}{u0}\to\subf{\alpha}{u1}$, then $u1$ has
the same \kl{polarity} as $u$ and $u0$ the opposite one. \AP For each
\kl{compound occurrence} $u$ we choose a new variable $\intro*\ev{u}$,
called an \intro{extension variable}, so that distinct occurrences receive
distinct variables and no $\ev{u}$ occurs in $\alpha$. For every
occurrence $u$ let
\[
\intro*\rv{u}=\begin{cases}\subf{\alpha}{u},&\text{if $u$ is a \kl{leaf}},\\
\ev{u},&\text{if $u$ is a \kl{compound occurrence}.}\end{cases}
\]
\AP For every \kl{compound occurrence} $u$ we call
$\intro*\body{u}=\rv{u0}\to\rv{u1}$ the \intro{body} of $u$, and we define
the \intro{defining premise} of $u$ by
\[
\intro*\dprem{u}=\begin{cases}\body{u}\to\ev{u},&\text{if $u$ is \kl{positive}},\\
\ev{u}\to\body{u},&\text{if $u$ is \kl{negative}.}\end{cases}
\]

\begin{definition}\label{def:S}
\AP Let $u_1,\dots,u_m$ be the \kl{compound occurrences} of $\alpha$,
listed in post-order; that is, $u_i$ precedes $u_j$ if and only if $u_i$
is \kl{strictly below} $u_j$, or $u_i=w0u'$ and $u_j=w1u''$ for some
words $w,u',u''$. The \intro{translation} of $\alpha$ is the formula
\[
\intro*\St(\alpha)=\dprem{u_1}\to\dots\to\dprem{u_m}\to\rv{\rt} .
\]
\end{definition}

If $\alpha$ is a variable, then $m=0$ and $\St(\alpha)=\alpha$. \AP Let $\intro*\decsub{\alpha}$ be the substitution that maps each
$\ev{u}$ to $\subf{\alpha}{u}$ and fixes all other variables; then
$\decsub{\alpha}(\rv{v})=\subf{\alpha}{v}$ for every occurrence $v$ of
$\alpha$.


In the following examples, each implication node of the syntax tree of
$\alpha$, that is, each \kl{compound occurrence} $u$, is annotated with
$\ev{u}$ and the \kl{polarity} of $u$; the left and right children of $u$
are $u0$ and $u1$.

\begin{example}\label{ex:peirce}
For Peirce's law $\alpha=((a\to b)\to a)\to a$:
\[
\begin{tikzpicture}[syntree,
  level 1/.style={sibling distance=26mm},
  level 2/.style={sibling distance=16mm},
  level 3/.style={sibling distance=9mm}]
\node[ann={$\ev{\rt}^{+}$}] {$\to$}
  child {node[ann={$\ev{0}^{-}$}] {$\to$}
    child {node[ann={$\ev{00}^{+}$}] {$\to$} child {node {$a$}} child {node {$b$}}}
    child {node {$a$}}}
  child {node {$a$}};
\end{tikzpicture}
\qquad
\begin{aligned}
\dprem{00}&=\body{00}\to\ev{00}=(a\to b)\to\ev{00},\\
\dprem{0}&=\ev{0}\to\body{0}=\ev{0}\to(\ev{00}\to a),\\
\dprem{\rt}&=\body{\rt}\to\ev{\rt}=(\ev{0}\to a)\to\ev{\rt}.
\end{aligned}
\]
Hence
\[
\St(\alpha)=((a\to b)\to\ev{00})\to(\ev{0}\to\ev{00}\to a)\to((\ev{0}\to a)\to\ev{\rt})\to\ev{\rt}.
\]
\end{example}

\begin{example}\label{ex:lc}
For the axiom $\alpha=((a\to b)\to c)\to((b\to a)\to c)\to c$ of
Gödel--Dummett logic:
\[
\begin{tikzpicture}[syntree,
  level 1/.style={sibling distance=38mm},
  level 2/.style={sibling distance=19mm},
  level 3/.style={sibling distance=14mm},
  level 4/.style={sibling distance=9mm}]
\node[ann={$\ev{\rt}^{+}$}] {$\to$}
  child {node[ann={$\ev{0}^{-}$}] {$\to$}
    child {node[ann={$\ev{00}^{+}$}] {$\to$} child {node {$a$}} child {node {$b$}}}
    child {node {$c$}}}
  child {node[ann={$\ev{1}^{+}$}] {$\to$}
    child {node[ann={$\ev{10}^{-}$}] {$\to$}
      child {node[ann={$\ev{100}^{+}$}] {$\to$} child {node {$b$}} child {node {$a$}}}
      child {node {$c$}}}
    child {node {$c$}}};
\end{tikzpicture}
\qquad
\begin{aligned}
\dprem{00}&=(a\to b)\to\ev{00},\\
\dprem{0}&=\ev{0}\to(\ev{00}\to c),\\
\dprem{100}&=(b\to a)\to\ev{100},\\
\dprem{10}&=\ev{10}\to(\ev{100}\to c),\\
\dprem{1}&=(\ev{10}\to c)\to\ev{1},\\
\dprem{\rt}&=(\ev{0}\to\ev{1})\to\ev{\rt}.
\end{aligned}
\]
Listing the \kl{defining premises} in this order (post-order), we obtain
\[
\begin{aligned}
\St(\alpha)={}&((a\to b)\to\ev{00})\to(\ev{0}\to\ev{00}\to c)
\to((b\to a)\to\ev{100})\to(\ev{10}\to\ev{100}\to c)\\
&\to((\ev{10}\to c)\to\ev{1})\to((\ev{0}\to\ev{1})\to\ev{\rt})\to\ev{\rt}.
\end{aligned}
\]
\end{example}
